% exemplo_uso_breqn.tex
% Este arquivo demonstra o uso do pacote breqn para quebra automática de equações

% Para testar, copie este conteúdo para uma seção do documento principal

\subsection{Exemplo de Uso do breqn}

\subsubsection{Equação Curta - Use equation Normal}

Para equações curtas, continue usando o ambiente tradicional:

\begin{equation}
E = mc^2
\label{eq:einstein_exemplo}
\end{equation}

\subsubsection{Equação Longa - Use dmath}

Para equações longas que podem ultrapassar a margem, use o ambiente \texttt{dmath}:

\begin{dmath}
\gamma(T,P) = \frac{1}{2A} \left[ G_{\text{slab}}(T,P,N_{\text{M}},N_{\text{W}},N_{\text{O}}) - N_{\text{M}}\mu_{\text{M}}(T,P) - N_{\text{W}}[g^{bulk}_{M_xWO_4}(T,P) - x\mu_{M}(T,P) - 4\mu_{O}(T,P)] - N_{\text{O}}\mu_{\text{O}}(T,P) \right]
\label{eq:gibbs_longa_exemplo}
\end{dmath}

A equação acima será automaticamente quebrada em várias linhas se não couber na largura da página.

\subsubsection{Múltiplas Equações com Quebra}

Use o ambiente \texttt{dgroup} para grupos de equações:

\begin{dgroup}
\begin{dmath}
\mu_{W} = g^{bulk}_{M_xWO_4}(T,P) - x\mu_{M}(T,P) - 4\mu_{O}(T,P)
\label{eq:mu_w_exemplo}
\end{dmath}
\begin{dmath}
\mu^{bulk}_{M_xWO_4}(T,P) = x\mu_{M}(T,P) + \mu_{W}(T,P) + 4\mu_{O}(T,P) \approx g^{bulk}_{M_xWO_4}(T,P)
\label{eq:mu_bulk_exemplo}
\end{dmath}
\end{dgroup}

\subsubsection{Referenciando Equações}

As equações podem ser referenciadas normalmente:
\begin{itemize}
    \item A Equação~\ref{eq:einstein_exemplo} mostra a famosa relação energia-massa.
    \item A Equação~\ref{eq:gibbs_longa_exemplo} descreve a energia livre de Gibbs da superfície.
    \item As Equações~\ref{eq:mu_w_exemplo} e~\ref{eq:mu_bulk_exemplo} relacionam potenciais químicos.
\end{itemize}

\subsubsection{Dicas de Uso}

\begin{enumerate}
    \item Use \texttt{dmath} quando suspeitar que a equação pode ser longa
    \item Use \texttt{dmath*} se não quiser numeração
    \item Mantenha \texttt{equation} para equações curtas e simples
    \item O \texttt{breqn} quebra automaticamente nos operadores (+, -, =, etc.)
    \item A indentação das linhas seguintes é automática
\end{enumerate}
